% =====================================================================
%  ANEXOS
% =====================================================================
\appendix
\addto\captionsspanish{\renewcommand{\appendixname}{Anexo}}
\renewcommand{\appendixname}{Anexo}

% ---------------------------------------------------------------------
\chapter{Panel de acciones}
\label{anx:panel}

Las 47 acciones del panel, ordenadas por sector. «Merval» indica si integra el S\&P~Merval;
«Volumen» es el porcentaje de ruedas con operaciones en los tres años; «Volatilidad» es la
volatilidad realizada promedio entre el 2/1/2025 y el 30/4/2026; «Ruedas con noticias» cuenta
las ruedas de la unión en que al menos una oración la menciona; «Tono propio» indica si tiene
noticias suficientes (90 ruedas o más) para medirlo. Se citan en la sección~\ref{sec:unidades}.

{\footnotesize\setlength{\tabcolsep}{3.5pt}
\begin{longtable}{@{}lL{3.2cm}L{2.35cm}C{1.1cm}C{1.3cm}C{1.75cm}C{1.4cm}C{1.1cm}@{}}
\toprule
\textbf{Ticker} & \textbf{Empresa} & \textbf{Sector} & \textbf{Merval} & \textbf{Volumen (\%)} & \textbf{Volatilidad (\%)} & \textbf{Ruedas con noticias} & \textbf{Tono propio} \\
\midrule
\endhead
\bottomrule
\caption{Panel de 47 acciones. Fuente: elaboración propia
(\carpeta{Lado Bursatil/01_scraping_acciones}, \carpeta{05_volatilidad_ranking} y
\carpeta{Unión/03_sentimiento_por_empresa}).}
\label{tab:panel}
\endlastfoot
CVH & Cablevision Holding & Comunicaciones & Sí & 100,0 & 39,3 & 8 & No \\
TECO2 & Telecom Argentina & Comunicaciones & Sí & 100,0 & 52,2 & 52 & No \\
COME & Sociedad Comercial del Plata & Conglomerado & Sí & 100,0 & 68,3 & 74 & No \\
CRES & Cresud & Conglomerado & Sí & 100,0 & 41,3 & 163 & Sí \\
GARO & Garovaglio y Zorraquín & Conglomerado & No & 97,0 & 70,5 & 1 & No \\
GRIM & Grimoldi & Consumo & No & 96,9 & 54,8 & 20 & No \\
HAVA & Havanna Holding & Consumo & No & 100,0 & 38,7 & 29 & No \\
LEDE & Ledesma & Consumo & No & 100,0 & 41,6 & 11 & No \\
LONG & Longvie & Consumo & No & 100,0 & 47,2 & 0 & No \\
MIRG & Mirgor & Consumo & Sí & 100,0 & 36,2 & 27 & No \\
MOLI & Molinos Río de la Plata & Consumo & No & 100,0 & 58,5 & 13 & No \\
MORI & Morixe Hermanos & Consumo & No & 100,0 & 62,9 & 7 & No \\
SAMI & San Miguel & Consumo & No & 100,0 & 64,7 & 2 & No \\
SEMI & Molinos Juan Semino & Consumo & No & 100,0 & 55,2 & 1 & No \\
CAPX & CAPEX SA & Energía & No & 100,0 & 52,8 & 0 & No \\
TGNO4 & Transportadora de Gas del Norte & Energía & Sí & 100,0 & 61,7 & 75 & No \\
TGSU2 & Transportadora de Gas del Sur & Energía & Sí & 100,0 & 53,8 & 189 & Sí \\
YPFD & YPF SA & Energía & Sí & 100,0 & 44,9 & 237 & Sí \\
BBAR & BBVA Banco Francés & Financiero & Sí & 100,0 & 65,9 & 229 & Sí \\
BHIP & Banco Hipotecario & Financiero & No & 100,0 & 66,4 & 47 & No \\
BMA & Banco Macro & Financiero & Sí & 100,0 & 60,9 & 184 & Sí \\
BPAT & Banco Patagonia & Financiero & No & 100,0 & 60,9 & 56 & No \\
BYMA & Bolsas y Mercados Argentinos & Financiero & Sí & 100,0 & 46,7 & 141 & Sí \\
GGAL & Grupo Financiero Galicia & Financiero & Sí & 100,0 & 55,8 & 195 & Sí \\
SUPV & Supervielle & Financiero & Sí & 100,0 & 74,4 & 223 & Sí \\
VALO & Banco de Valores & Financiero & Sí & 100,0 & 40,2 & 50 & No \\
AGRO & Agrometal & Industrial & No & 100,0 & 64,3 & 2 & No \\
AUSO & Autopistas del Sol & Industrial & No & 100,0 & 55,0 & 7 & No \\
OEST & Grupo Concesionario del Oeste & Industrial & No & 99,7 & 46,9 & 1 & No \\
ALUA & Aluar & Materiales & Sí & 100,0 & 50,2 & 94 & Sí \\
CELU & Celulosa Argentina & Materiales & No & 100,0 & 104,6 & 13 & No \\
FIPL & Fiplasto & Materiales & No & 99,9 & 41,3 & 0 & No \\
LOMA & Loma Negra & Materiales & No & 100,0 & 51,2 & 175 & Sí \\
RIGO & Rigolleau & Materiales & No & 95,8 & 58,1 & 4 & No \\
TXAR & Ternium Argentina & Materiales & Sí & 100,0 & 48,1 & 105 & Sí \\
CTIO & Consultatio & Real Estate & No & 100,0 & 37,5 & 38 & No \\
GCDI & GCDI S.A. (ex TGLT) & Real Estate & No & 94,9 & 46,2 & 1 & No \\
IRSA & IRSA Inversiones y Representaciones & Real Estate & No & 100,0 & 43,7 & 186 & Sí \\
BOLT & Boldt & Tecnología & No & 99,7 & 44,1 & 1 & No \\
CEPU & Central Puerto & Utilities & Sí & 100,0 & 53,7 & 193 & Sí \\
CGPA2 & Camuzzi Gas Pampeana & Utilities & No & 100,0 & 67,5 & 21 & No \\
DGCU2 & Distribuidora de Gas Cuyana & Utilities & No & 100,0 & 55,2 & 5 & No \\
EDN & Edenor & Utilities & Sí & 100,0 & 65,5 & 213 & Sí \\
GBAN & Naturgy BAN & Utilities & No & 100,0 & 58,0 & 4 & No \\
METR & MetroGas & Utilities & No & 100,0 & 73,6 & 112 & Sí \\
PAMP & Pampa Energía & Utilities & Sí & 100,0 & 44,5 & 164 & Sí \\
TRAN & Transener & Utilities & Sí & 100,0 & 65,9 & 116 & Sí \\
\end{longtable}
}

% ---------------------------------------------------------------------
\chapter{Diccionario de variables}
\label{anx:diccionario}

Columnas de los dos datasets de la unión: \var{dataset\_mercado.csv} (una fila por rueda) y
\var{dataset\_panel.csv} (una fila por acción y rueda). «Tabla» indica en cuál está cada
variable; «ambos», que está en las dos. El archivo original es
\carpeta{Unión/04_dataset_unificado/diccionario_variables.csv}.

{\footnotesize\setlength{\tabcolsep}{3.5pt}
\begin{longtable}{@{}L{3.0cm}L{1.2cm}L{5.1cm}L{1.9cm}L{3.0cm}@{}}
\toprule
\textbf{Variable} & \textbf{Tabla} & \textbf{Qué es} & \textbf{Unidad} & \textbf{De dónde sale} \\
\midrule
\endhead
\bottomrule
\caption{Diccionario de variables de los datasets unificados. Fuente: elaboración propia.}
\label{tab:diccionario}
\endlastfoot
\nolinkurl{rueda} & ambos & Día de bolsa (BYMA) & fecha & Lado \nolinkurl{Bursatil/03_vix_argentino} (calendario de \^{}MERV) \\
\nolinkurl{merv_cierre} & mercado & Cierre del índice Merval & puntos & Lado \nolinkurl{Bursatil/03_vix_argentino} \\
\nolinkurl{ret_merval} & ambos & Retorno diario del Merval & proporción (0,01 = 1\%) & Lado \nolinkurl{Bursatil/03_vix_argentino} \\
\nolinkurl{vix_arg} & ambos & VIX argentino: volatilidad realizada de 20 ruedas, anualizada & \% & Lado \nolinkurl{Bursatil/03_vix_argentino} \\
\nolinkurl{regimen} & ambos & Régimen de mercado (HMM, probabilidades suavizadas): Calmo, Moderado o Tenso & categoría & Lado \nolinkurl{Bursatil/06_hmm_regimen_mercado} \\
\nolinkurl{regimen_filtrado} & ambos & Régimen con probabilidades filtradas (solo usa información hasta ese día) & categoría & Lado \nolinkurl{Bursatil/06_hmm_regimen_mercado} \\
\nolinkurl{prob_tenso} & ambos & Probabilidad suavizada del régimen Tenso & 0 a 1 & Lado \nolinkurl{Bursatil/06_hmm_regimen_mercado} \\
\nolinkurl{prob_tenso_filtrada} & mercado & Probabilidad filtrada del régimen Tenso & 0 a 1 & Lado \nolinkurl{Bursatil/06_hmm_regimen_mercado} \\
\nolinkurl{monto_operado} & mercado & Cierre $\times$ volumen, sumado sobre las 47 acciones & pesos & Lado \nolinkurl{Bursatil/01_scraping_acciones} (U4) \\
\nolinkurl{monto_anormal} & mercado & log(monto) $-$ promedio de log(monto) de las 20 ruedas previas & log (0 = habitual) & U4 \\
\nolinkurl{n_notas} & ambos & Notas que recibe la rueda (regla de las 17:00) & notas & U1 \\
\nolinkurl{dias_cubiertos} & mercado & Días de calendario que junta la rueda (1 normal, 3 un lunes, 4+ tras feriado) & días & U2 \\
\nolinkurl{atencion_anormal} & ambos & Atención de los medios: notas comparadas con lo habitual para ese tipo de rueda y las 20 previas & log (0 = habitual) & U2 \\
\nolinkurl{sent_neto} & ambos & Tono de las noticias: promedio del sentimiento neto de las notas de la rueda & puntos ($-$100 a +100) & U2 (sentimiento de Gino) \\
\nolinkurl{sent_neto_portal} & ambos & Tono ajustado por portal (a cada nota se le resta el promedio de su portal) & puntos & U2 \\
\nolinkurl{sent_anormal} & ambos & Tono anormal: tono de la rueda $-$ promedio del tono de las 20 ruedas previas. Variable de tono del VAR & puntos (0 = habitual) & U4 (decisión de U5) \\
\nolinkurl{sent_portal_anormal} & ambos & Lo mismo con el tono ajustado por portal & puntos (0 = habitual) & U4 (decisión de U5) \\
\nolinkurl{pct_pos, pct_neg, pct_neu, pct_mixto} & mercado & Porcentaje promedio de oraciones de cada clase en las notas de la rueda & \% & U2 \\
\nolinkurl{z_shock} & mercado & Distancia del tono a lo habitual, en veces el error esperable por la cantidad de notas & veces & U2 \\
\nolinkurl{shock, shock_pos, shock_neg} & ambos & Rueda de tono extremo (|z| > 2,5), total, positiva y negativa & 0/1 & U2 \\
\nolinkurl{evento_id} & ambos & Número de evento (shocks del mismo signo en ruedas seguidas = un evento); 0 si no hay & número & U2 \\
\nolinkurl{dia_electoral} & ambos & Rueda posterior a una elección (08/09/2025 y 27/10/2025) & texto & U2 \\
\nolinkurl{eco_electoral} & ambos & Rueda dentro de las 20 posteriores a un día electoral (eco en el VIX y el régimen) & 0/1 & U4 \\
\nolinkurl{abril_2026} & ambos & Rueda de abril de 2026 (hueco de iProfesional en el corpus) & 0/1 & U4 \\
\nolinkurl{ticker, empresa, sector, oficial} & panel & Identificación de la acción; oficial = está en el S\&P Merval & texto & Lado \nolinkurl{Bursatil/01_scraping_acciones} \\
\nolinkurl{grupo_baja_liquidez} & panel & Integra el grupo de 8 del clustering, explicado por baja liquidez & 0/1 & Lado \nolinkurl{Bursatil/02_clustering_acciones} \\
\nolinkurl{cuatro_iliquidas} & panel & Es una de las 4 acciones con más días sin operar (GCDI, RIGO, GRIM, GARO) & 0/1 & Lado \nolinkurl{Bursatil/02_clustering_acciones} \\
\nolinkurl{cierre} & panel & Precio de cierre ajustado & pesos & Lado \nolinkurl{Bursatil/01_scraping_acciones} \\
\nolinkurl{retorno} & panel & Retorno diario de la acción (pct\_change sobre la serie completa) & proporción & U4 (misma definición que la carpeta 04) \\
\nolinkurl{volumen} & panel & Acciones operadas en la rueda & acciones & Lado \nolinkurl{Bursatil/01_scraping_acciones} \\
\nolinkurl{monto} & panel & Cierre $\times$ volumen & pesos & U4 \\
\nolinkurl{vol_anormal} & panel & log(1 + volumen) $-$ promedio de log(1 + volumen) de las 20 ruedas previas & log (0 = habitual) & U4 \\
\nolinkurl{sin_operar} & panel & La acción no operó (volumen 0) & 0/1 & U4 \\
\nolinkurl{precio_quieto} & panel & El cierre es igual al de la rueda anterior & 0/1 & U4 \\
\nolinkurl{hubo_noticia} & panel & Al menos una oración de las notas de la rueda menciona a la empresa & 0/1 & U3 \\
\nolinkurl{sentimiento_propio} & panel & La empresa tiene noticias en al menos 90 ruedas y usa su propio sentimiento & 0/1 & U3 \\
\nolinkurl{n_articulos_empresa, n_oraciones_empresa} & panel & Notas y oraciones de la rueda que mencionan a la empresa & cantidad & U3 \\
\nolinkurl{sent_neto_empresa} & panel & Tono de las oraciones que mencionan a la empresa; vacío si no hubo noticia & puntos & U3 (cruce de Gino) \\
\nolinkurl{peso_tono_propio} & panel & Cuánto se confía en las oraciones de la empresa: n / (n + K), con K estimado de los datos & 0 a 1 & U4 (decisión de U6) \\
\nolinkurl{tono_propio} & panel & Tono propio: peso $\times$ (tono de la empresa $-$ tono del mercado); 0 si no hubo noticia. Variable de empresa del VAR & puntos (0 = como el mercado) & U4 (decisión de U6) \\
\end{longtable}
}

% ---------------------------------------------------------------------
\chapter{Alineación de las notas con las ruedas}
\label{anx:alineacion}

\textbf{Horas sospechosas.} Se marcó como sospechosa toda hora que concentra más del 3\% de
las notas de un portal en un mismo minuto. Las detectadas son: Ámbito 00:00 y 07:00; Clarín
04:00, 06:01, 07:00, 07:05, 07:08 y 09:00; iProfesional 07:00; La Nación 01:05 y 02:00; TN
05:44, 05:48 y 05:55. Suman 1.881 notas (13,4\%). La mayoría cae de madrugada, por lo que la
regla las asigna a la rueda del mismo día; si en realidad salieron la noche anterior ya están
en la rueda correcta, y si salieron después de las 17:00 deberían ir a la siguiente. No hay
forma de saberlo con los datos, y por eso los resultados se repitieron sin ellas.

\textbf{Esperas.} La espera más larga entre la publicación y la rueda es de cinco días: una nota
del viernes 13/02/2026 a la noche cuenta para el miércoles 18/02/2026, porque el lunes y el
martes fueron feriados de Carnaval.

\textbf{Controles automáticos (10 de 10).} Cada nota tiene una rueda o un motivo para no
tenerla; ninguna va a una rueda anterior a su publicación; ninguna publicada tras el cierre
queda en la rueda del mismo día; las URL y los portales coinciden uno a uno con el archivo de
sentimiento del lado de noticias; tomando la fecha sin convertir se reproduce la fecha de su
serie diaria; todas las ruedas tienen noticias; y la cantidad de ruedas, el mínimo y la mediana
de notas por rueda y las notas con hora sospechosa coinciden con el diagnóstico previo.

\textbf{Ruedas de tono extremo.} Las 19 ruedas marcadas fueron: 18/06, 15/09, 24/09, 25/09,
13/10, 16/10, 27/10, 30/10, 03/11, 14/11, 25/11, 01/12 y 23/12/2025; 23/01, 26/01, 04/02,
13/02, 16/03 y 31/03/2026. Siete coinciden con un shock de la serie diaria del lado de
noticias; las diferencias se deben a la regla de la rueda, al ajuste por portal y a que la
regla por muestreo no marca días con pocas notas.

% ---------------------------------------------------------------------
\chapter{Detalle de las pruebas de Granger}
\label{anx:granger}

$p$ corregidos por Benjamini-Hochberg dentro de cada nivel. ${}^{\bullet}$ significativo
después de la corrección; ${}^{\circ}$ $p<0{,}05$ solo sin corregir. «Rez.» es la cantidad de
rezagos elegida por el criterio de Akaike. Se comentan en la sección~\ref{sec:res-u6}.

{\small
\begin{longtable}{@{}lcC{2.55cm}C{2.55cm}C{2.55cm}C{2.55cm}@{}}
\toprule
\textbf{Empresa} & \textbf{Rez.} & \textbf{Propio \flecha{} retorno} & \textbf{Propio \flecha{} volumen} & \textbf{Retorno \flecha{} propio} & \textbf{Volumen \flecha{} propio} \\
\midrule
\endhead
\bottomrule
\caption{Granger con el tono propio de las 17 empresas. Fuente: elaboración propia
(\carpeta{Unión/06_var_granger/granger_resultados.csv}).}
\label{tab:granger-propio}
\endlastfoot
ALUA & 1 & 0,6272 & 0,9768 & 0,0593$^{\circ}$ & 0,4127 \\
BBAR & 1 & 0,5969 & 0,9001 & $<$0,0001$^{\bullet}$ & 0,8505 \\
BMA & 1 & 0,7880 & 0,1776 & 0,0091$^{\bullet}$ & 0,9206 \\
BYMA & 1 & 0,9768 & 0,9206 & $<$0,0001$^{\bullet}$ & 0,5643 \\
CEPU & 1 & 0,6321 & 0,3490 & $<$0,0001$^{\bullet}$ & 0,9565 \\
CRES & 1 & 0,4127 & 0,3169 & 0,0006$^{\bullet}$ & 0,0584$^{\circ}$ \\
EDN & 1 & 0,8505 & 0,1442$^{\circ}$ & $<$0,0001$^{\bullet}$ & 0,8505 \\
GGAL & 1 & 0,5643 & 0,4541 & 0,0002$^{\bullet}$ & 0,9001 \\
IRSA & 1 & 0,5969 & 0,0215$^{\bullet}$ & 0,0002$^{\bullet}$ & 0,6272 \\
LOMA & 1 & 0,7883 & 0,3251 & $<$0,0001$^{\bullet}$ & 0,9450 \\
METR & 1 & 0,6271 & 0,3251 & $<$0,0001$^{\bullet}$ & 0,8505 \\
PAMP & 1 & 0,6272 & 0,9523 & 0,9206 & 0,5969 \\
SUPV & 1 & 0,5513 & 0,4127 & $<$0,0001$^{\bullet}$ & 0,9206 \\
TGSU2 & 1 & 0,9001 & 0,5969 & 0,0002$^{\bullet}$ & 0,5969 \\
TRAN & 1 & 0,3080 & 0,5969 & 0,0116$^{\bullet}$ & 0,8505 \\
TXAR & 1 & 0,6623 & 0,3251 & 0,0263$^{\bullet}$ & 0,9206 \\
YPFD & 1 & 0,9351 & 0,9206 & 0,0120$^{\bullet}$ & 0,3251 \\
\end{longtable}

\begin{longtable}{@{}lcC{2.55cm}C{2.55cm}C{2.55cm}C{2.55cm}@{}}
\toprule
\textbf{Empresa} & \textbf{Rez.} & \textbf{Tono \flecha{} retorno} & \textbf{Tono \flecha{} volumen} & \textbf{Retorno \flecha{} tono} & \textbf{Volumen \flecha{} tono} \\
\midrule
\endhead
\bottomrule
\caption{Granger de las 47 empresas con el tono anormal del mercado. Fuente: elaboración
propia (\carpeta{Unión/06_var_granger/granger_resultados.csv}).}
\label{tab:granger-empresas}
\endlastfoot
AGRO & 2 & 0,4426 & 0,5936 & 0,0006$^{\bullet}$ & 0,0471$^{\bullet}$ \\
ALUA & 1 & 0,7962 & 0,6495 & 0,1518 & 0,3506 \\
AUSO & 4 & 0,4426 & 0,2305 & 0,0024$^{\bullet}$ & 0,5449 \\
BBAR & 4 & 0,3119 & 0,1817 & 0,0023$^{\bullet}$ & 0,6873 \\
BHIP & 4 & 0,2069 & 0,3417 & $<$0,0001$^{\bullet}$ & 0,8127 \\
BMA & 1 & 0,2765 & 0,9943 & 0,0034$^{\bullet}$ & 0,5034 \\
BOLT & 2 & 0,1170$^{\circ}$ & 0,6640 & 0,1416$^{\circ}$ & 0,9375 \\
BPAT & 1 & 0,2347 & 0,3628 & 0,0003$^{\bullet}$ & 0,1054$^{\circ}$ \\
BYMA & 3 & 0,9188 & 0,1293$^{\circ}$ & 0,0038$^{\bullet}$ & 0,3691 \\
CAPX & 2 & 0,5481 & 0,6259 & 0,0003$^{\bullet}$ & 0,9607 \\
CELU & 1 & 0,7656 & 0,9672 & 0,7001 & 0,0633$^{\circ}$ \\
CEPU & 4 & 0,2086 & 0,3628 & 0,0003$^{\bullet}$ & 0,3628 \\
CGPA2 & 3 & 0,3691 & 0,4717 & $<$0,0001$^{\bullet}$ & 0,3119 \\
COME & 1 & 0,2847 & 0,5446 & 0,0003$^{\bullet}$ & 0,1050$^{\circ}$ \\
CRES & 1 & 0,6495 & 0,6126 & 0,2107 & 0,6918 \\
CTIO & 2 & 0,4717 & 0,4426 & 0,0005$^{\bullet}$ & 0,8127 \\
CVH & 1 & 0,5875 & 0,1247$^{\circ}$ & 0,0042$^{\bullet}$ & 0,0236$^{\bullet}$ \\
DGCU2 & 1 & 0,4252 & 0,7780 & $<$0,0001$^{\bullet}$ & 0,5936 \\
EDN & 1 & 0,5875 & 0,4299 & 0,0003$^{\bullet}$ & 0,7448 \\
FIPL & 1 & 0,8180 & 0,5707 & 0,0024$^{\bullet}$ & 0,2107 \\
GARO & 1 & 0,6866 & 0,9502 & 0,0005$^{\bullet}$ & 0,7100 \\
GBAN & 1 & 0,3628 & 0,2181 & $<$0,0001$^{\bullet}$ & 0,3249 \\
GCDI & 1 & 0,4808 & 0,0503$^{\circ}$ & 0,0006$^{\bullet}$ & 0,2847 \\
GGAL & 1 & 0,2765 & 0,2472 & 0,0005$^{\bullet}$ & 0,1424$^{\circ}$ \\
GRIM & 1 & 0,9271 & 0,1538 & 0,1416$^{\circ}$ & 0,6373 \\
HAVA & 1 & 0,8926 & 0,1817 & 0,0784$^{\circ}$ & 0,5160 \\
IRSA & 1 & 0,6905 & 0,1424$^{\circ}$ & 0,3628 & 0,6332 \\
LEDE & 1 & 0,8343 & 0,3119 & 0,0471$^{\bullet}$ & 0,4426 \\
LOMA & 2 & 0,5936 & 0,6495 & 0,0135$^{\bullet}$ & 0,4808 \\
LONG & 1 & 0,6495 & 0,0126$^{\bullet}$ & 0,0030$^{\bullet}$ & 0,0024$^{\bullet}$ \\
METR & 4 & 0,4474 & 0,2107 & $<$0,0001$^{\bullet}$ & 0,9769 \\
MIRG & 2 & 0,4252 & 0,4808 & 0,0017$^{\bullet}$ & 0,1538 \\
MOLI & 1 & 0,4717 & 0,4424 & 0,0024$^{\bullet}$ & 0,0652$^{\circ}$ \\
MORI & 2 & 0,2750 & 0,5875 & 0,0462$^{\bullet}$ & 0,4717 \\
OEST & 1 & 0,6545 & 0,8127 & 0,0006$^{\bullet}$ & 0,4426 \\
PAMP & 1 & 0,3964 & 0,1538 & 0,0024$^{\bullet}$ & 0,4717 \\
RIGO & 1 & 0,2157 & 0,5136 & 0,0185$^{\bullet}$ & 0,3690 \\
SAMI & 2 & 0,6905 & 0,7260 & 0,0245$^{\bullet}$ & 0,1518 \\
SEMI & 3 & 0,7292 & 0,3691 & 0,0077$^{\bullet}$ & 0,4857 \\
SUPV & 1 & 0,3029 & 0,0498$^{\bullet}$ & 0,0030$^{\bullet}$ & 0,4426 \\
TECO2 & 1 & 0,9672 & 0,7001 & 0,0327$^{\bullet}$ & 0,1424$^{\circ}$ \\
TGNO4 & 1 & 0,6332 & 0,8127 & $<$0,0001$^{\bullet}$ & 0,6312 \\
TGSU2 & 1 & 0,2654 & 0,5936 & 0,0154$^{\bullet}$ & 0,1240$^{\circ}$ \\
TRAN & 1 & 0,7962 & 0,6980 & $<$0,0001$^{\bullet}$ & 0,2750 \\
TXAR & 1 & 0,7100 & 0,2847 & 0,0784$^{\circ}$ & 0,1817 \\
VALO & 1 & 0,9849 & 0,3628 & 0,0044$^{\bullet}$ & 0,1293$^{\circ}$ \\
YPFD & 1 & 0,4717 & 0,1538 & 0,0044$^{\bullet}$ & 0,3417 \\
\end{longtable}
}

\begin{table}[H]
\centering\footnotesize
\begin{tabular}{@{}lcccc@{}}
\toprule
\textbf{Sector (acciones)} & \textbf{Tono \flecha{} retorno} & \textbf{Tono \flecha{} volumen} & \textbf{Retorno \flecha{} tono} & \textbf{Volumen \flecha{} tono} \\
\midrule
Comunicaciones (2) & — & — & Sí (+) & Sí (+) \\
Conglomerado (3) & — & Sí ($-$) & Sí (+) & — \\
Consumo (9) & — & sin corregir & Sí (+) & Sí (+) \\
Energía (4) & — & — & Sí (+) & — \\
Financiero (8) & — & — & Sí (+) & — \\
Industrial (3) & — & — & Sí (+) & — \\
Materiales (6) & — & — & Sí (+) & — \\
Real Estate (3) & — & sin corregir & Sí (+) & — \\
Tecnología (1) & sin corregir & — & sin corregir & — \\
Utilities (8) & — & — & Sí (+) & — \\
\bottomrule
\end{tabular}
\caption{Granger por sector. «Sí» indica significativo después de la corrección, con el signo
del efecto; «sin corregir», $p<0{,}05$ solo sin corrección. Fuente: elaboración propia
(\carpeta{Unión/06_var_granger}).}
\label{tab:granger-sectores}
\end{table}

% ---------------------------------------------------------------------
\chapter{Especificaciones escritas antes de los resultados}
\label{anx:especificaciones}

Las dos pruebas más directas de la hipótesis se especificaron por escrito y se guardaron antes
de correr ningún cálculo. Si algo hubiera que cambiar, se agrega como «cambio posterior», con la
fecha y el motivo, y se informan las dos versiones. Los textos completos están en
\carpeta{Unión/06b_noticias_preapertura/especificacion_previa.md} y
\carpeta{Unión/07b_shock_volumen_noticias/especificacion_previa.md}.

\section{Notas publicadas con el mercado cerrado (U6b)}

\begin{itemize}
  \item \textbf{Pregunta.} ¿El tono de las noticias publicadas con el mercado cerrado anticipa
        el precio de la rueda siguiente? Esas notas salen antes de cualquier operación de la
        rueda, así que la rueda no puede haberlas causado.
  \item \textbf{Notas.} Las asignadas a la rueda $t$ que no se publicaron entre las 11:00 y las
        17:00 del día $t$, sin las de hora sospechosa.
  \item \textbf{Variables.} Tono pre-apertura anormal del mercado (menos el promedio de las 20
        ruedas previas); tono propio pre-apertura de las 17 empresas, con la misma fórmula del
        tono propio y el mismo $K$; retorno de apertura (apertura de $t$ / cierre de $t-1$
        $-$ 1) y de la sesión (cierre de $t$ / apertura de $t$ $-$ 1). El mercado es el
        promedio simple de las 47 acciones.
  \item \textbf{Modelos.} Mínimos cuadrados con errores robustos HC3. Mercado: retorno $= a +
        b\cdot$tono pre-apertura anormal $+ c\cdot$retorno del día anterior $+ d\cdot$rueda tras
        día sin mercado. Empresas: panel con efecto fijo por empresa, tono propio y tono del
        mercado, retorno propio del día anterior y errores agrupados por rueda; y el mismo
        modelo para cada empresa por separado.
  \item \textbf{Evidencia.} La hipótesis se apoya si $b>0$ y es significativo ($p<0{,}05$) en
        el retorno de apertura, en el mercado o en el panel, con corrección de
        Benjamini-Hochberg dentro de cada grupo (4 pruebas principales y 34 por empresa). Se
        controla la calibración con el tono mezclado al azar.
  \item \textbf{Robustez agregada antes de correr.} Sin los dos días electorales.
\end{itemize}

\section{Shock de cantidad de noticias (U7b)}

\begin{itemize}
  \item \textbf{Shock.} Puntaje $z$ de la cantidad de la rueda contra las 30 ruedas previas
        (al menos 20 con dato), umbral $z>2$. Mercado: notas de la rueda. Empresas (las 17 con
        tono propio): artículos que la mencionan, con al menos 3 en la rueda. Ruedas seguidas
        con shock cuentan como un solo evento.
  \item \textbf{Medición.} Igual que el estudio de eventos de tono: ventanas $[-25,-4]$ y
        $[-3,+5]$; media constante para el S\&P~Merval y modelo de mercado para cada acción; días
        electorales fuera de la estimación.
  \item \textbf{Medidas.} $|$retorno anormal$|$ del día $[0]$; volumen anormal antes, el día y
        después; CAR antes, el día y después.
  \item \textbf{Significancia.} Contra 5.000 sorteos de fechas al azar, con Benjamini-Hochberg
        sobre las 14 pruebas principales (2 niveles $\times$ 7 medidas).
  \item \textbf{Evidencia.} H1 se apoya si, después de la corrección, el $|$retorno anormal$|$
        del día, el volumen anormal del día o de después, o el CAR de después son
        significativos. Si solo aparece movimiento «antes», es mercado \flecha{} noticias.
  \item \textbf{Robustez fijada de antemano.} Umbral 2,5; $z$ sobre $\log(1+\text{cantidad})$;
        con los eventos de ventana electoral.
  \item \textbf{Cambio posterior 1} (después de ver la versión original). Los 13 eventos del
        mercado caían en la primera rueda tras un día sin mercado, que junta las notas de tres o
        cuatro días. La cantidad pasó a medirse por día de calendario; todo lo demás quedó
        igual y se informan las dos versiones.
  \item \textbf{Agregado posterior 2} (antes de correrlo). Prueba pre-apertura: solo notas
        publicadas con el mercado cerrado, con el salto de apertura anormal como medida central
        y Benjamini-Hochberg sobre 6 pruebas.
  \item \textbf{Agregado posterior 3} (descriptivo). Sensibilidad del salto de apertura del
        mercado sacando un evento por vez.
\end{itemize}

% ---------------------------------------------------------------------
\chapter{Reproducibilidad y organización del proyecto}
\label{anx:repro}

El proyecto está en la carpeta \carpeta{tesis version 3.0}, con tres partes que no se
modifican entre sí. Cada carpeta tiene su script, sus resultados en CSV, sus controles
automáticos y un PDF que la documenta.

{\small
\begin{longtable}{@{}L{5.6cm}L{8.2cm}@{}}
\toprule
\textbf{Carpeta} & \textbf{Qué contiene} \\
\midrule
\endhead
\multicolumn{2}{@{}l}{\textit{Lado Bursatil}} \\
\carpeta{00_mapa_lado_bursatil} & Mapa del lado bursátil \\
\carpeta{01_scraping_acciones} & Selección del panel y precios diarios de las 47 acciones \\
\carpeta{02_clustering_acciones} & Clustering y análisis del grupo de baja liquidez \\
\carpeta{03_vix_argentino} & VIX argentino y ventana de fechas \\
\carpeta{04_adf} & Estacionariedad de precios y retornos \\
\carpeta{05_volatilidad_ranking} & Volatilidad por acción, ranking y picos propios \\
\carpeta{06_hmm_regimen_mercado} & Regímenes de mercado \\
\midrule
\multicolumn{2}{@{}l}{\textit{Lado Noticias}} \\
\carpeta{noticias} & Recolección, auditoría y depuración del corpus \\
\carpeta{sentimiento_finbert} & Etiquetado, FinBETO, medición de calidad y serie diaria \\
\carpeta{ner_entidades} & Reconocimiento de entidades y cruce con el sentimiento \\
\midrule
\multicolumn{2}{@{}l}{\textit{Unión}} \\
\carpeta{00_pipeline_union} & Pipeline, diagnóstico del cruce y decisiones previas \\
\carpeta{01_alineacion_noticias_ruedas} & U1: cada nota a su rueda (10 controles) \\
\carpeta{02_sentimiento_por_rueda} & U2: tono, atención y ruedas de tono extremo (7 controles) \\
\carpeta{03_sentimiento_por_empresa} & U3: tono de cada empresa por rueda (7 controles) \\
\carpeta{04_dataset_unificado} & U4: datasets y diccionario de variables (8 controles) \\
\carpeta{05_adf_union} & U5: estacionariedad (5 controles) \\
\carpeta{06_var_granger} & U6: VAR, Granger e impulso-respuesta (6 controles) \\
\carpeta{06b_noticias_preapertura} & U6b: notas con el mercado cerrado (6 controles) \\
\carpeta{07_event_study} & U7: estudio de eventos de tono extremo (5 controles) \\
\carpeta{07b_shock_volumen_noticias} & U7b: shock de cantidad de noticias (7 controles) \\
\carpeta{08_robustez} & U8: robustez (5 controles) \\
\carpeta{09_entrega_final} & Documento de la entrega, mapa de la tesis, plan actualizado y \texttt{correr\_todo.py} \\
\carpeta{10_papers_union} & Papers relacionados con cada prueba de la unión \\
\carpeta{11_papers_lado_noticias} & Papers relacionados con el corpus, el sentimiento y el NER \\
\bottomrule
\caption{Organización del proyecto. Fuente: elaboración propia.}
\label{tab:carpetas}
\end{longtable}
}

\textbf{Cómo se rehace la unión.} Desde \carpeta{Unión/09_entrega_final}, el comando
\texttt{python correr\_todo.py} corre los pasos U1 a U8 en orden (unos dos minutos y medio),
verifica que pasen los 66 controles automáticos y compara la huella de cada uno de los 64
archivos de resultados con la corrida anterior. Los documentos PDF se regeneran con los scripts
de \carpeta{99_generadores_documentacion}. Esta tesina se compila desde
\carpeta{Informe Final} con \texttt{python compilar\_tesina.py}, que copia las figuras, arma la
bibliografía a partir de los catálogos de papers y las tablas de los anexos a partir de los CSV.

% ---------------------------------------------------------------------
\chapter{Bibliografía por tema}
\label{anx:bib-tema}

Cada paso del trabajo tiene su propia carpeta de papers, con las referencias en
\texttt{referencias.bib} o \texttt{referencias\_papers.md}. Esta tabla agrupa las referencias de
la bibliografía según el tema al que respaldan; el archivo de cada una figura en la
Bibliografía.

{\small
\begin{longtable}{@{}L{4.2cm}L{9.6cm}@{}}
\toprule
\textbf{Tema} & \textbf{Referencias} \\
\midrule
\endhead
\bottomrule
\caption{Referencias agrupadas por tema. Fuente: elaboración propia.}
\label{tab:bib-tema}
\endlastfoot
\multicolumn{2}{@{}l}{\textit{Marco general}} \\
Medios y mercado & \citet{tetlock2007}; \citet{garcia2013}; \citet{fang2009}; \citet{engelberg2011}; \citet{dougal2012}; \citet{verma2026}; \citet{fama1970} \\
Metodología del proceso & \citet{wirth2000crispdm}; \citet{seabold2010} \\
\midrule
\multicolumn{2}{@{}l}{\textit{Lado bursátil (\carpeta{Lado Bursatil/*/papers})}} \\
00 Mapa & \citet{verma2026}; \citet{calandroni2026}; \citet{mukherjee2025}; \citet{amburgey2025} \\
01 Selección de acciones & \citet{nanyonga2025}; \citet{lesmond1999} \\
02 Clustering & \citet{mantegna1999}; \citet{rousseeuw1987}; \citet{onnela2003}; \citet{marti2021}; \citet{pharasi2020}; \citet{ghosh2025}; \citet{arratia2013}; \citet{ramirez2025}; \citet{alexandre2021}; \citet{huth2011}; \citet{busca2023}; \citet{lomackinlay1990} \\
03 VIX argentino & \citet{segovia2022}; \citet{segovia2022wp}; \citet{larre2019}; \citet{otuki2008}; \citet{calandroni2026}; \citet{sakowski2023}; \citet{papaioannou2014}; \citet{riskmetrics1996}; \citet{xenofontos2014} \\
04 Estacionariedad & \citet{dickey1979}; \citet{zablotsky2002}; \citet{worthington2003}; \citet{medina2013}; \citet{jain2020}; \citet{malafeyev2017}; \citet{kulikova2024}; \citet{rochafilho2019} \\
05 Volatilidad por acción & \citet{kalnina2024}; \citet{papaioannou2014}; \citet{otuki2008} \\
06 Régimen de mercado & \citet{hamilton1989}; \citet{schwarz1978}; \citet{kim1994}; \citet{ataurima2016}; \citet{caballero2016}; \citet{blake2025}; \citet{ataei2025}; \citet{majumder2023}; \citet{rabiner1989} \\
\midrule
\multicolumn{2}{@{}l}{\textit{Lado de noticias (\carpeta{Unión/11_papers_lado_noticias})}} \\
Recolección del corpus & \citet{barbaresi2021}; \citet{garg2025wayback}; \citet{foley2026wcxb}; \citet{hatwar2024author}; \citet{dodge2021c4}; \citet{lee2021dedup}; \citet{fronzetti2026gdelt} \\
Modelo de sentimiento & \citet{canete2020beto}; \citet{devlin2019bert}; \citet{araci2019finbert}; \citet{gururangan2020dapt}; \citet{malo2014phrasebank}; \citet{loughran2011}; \citet{perez2021pysentimiento}; \citet{gutierrez2022maria}; \citet{pan2023peerj}; \citet{marin2023finances}; \citet{garciadiaz2023ieee} \\
Etiquetado y evaluación & \citet{gilardi2023chatgpt}; \citet{aguda2024llmannot}; \citet{calderon2025alttest}; \citet{gemma2026gemma4}; \citet{kapoor2023leakage}; \citet{xue2026temporal} \\
Entidades nombradas & \citet{tjongkimsang2002conll}; \citet{honnibal2020spacy}; \citet{ujiie2021dict}; \citet{golde2025finerweb}; \citet{asgarieh2024salient} \\
\midrule
\multicolumn{2}{@{}l}{\textit{Unión (\carpeta{Unión/10_papers_union})}} \\
U2 Tono por rueda & \citet{brana2016}; \citet{rodriguezgarcia2023}; \citet{stan2026}; \citet{mavillonio2026}; \citet{carbo2026}; \citet{shiffler1988} \\
U3 Tono por empresa & \citet{tang2023finentity}; \citet{sinha2023sentfin}; \citet{chen2025topics} \\
U4 Atención y volumen & \citet{cao2025hype}; \citet{rojas2015}; \citet{efron1975}; \citet{fang2009}; \citet{da2011} \\
U5 Estacionariedad del tono & \citet{correarocha2026}; \citet{dickey1979}; \citet{kpss1992} \\
U6 VAR y Granger & \citet{granger1969}; \citet{sims1980}; \citet{lutkepohl2005}; \citet{bh1995}; \citet{verma2026}; \citet{mukherjee2025}; \citet{chniguir2022}; \citet{wu2026}; \citet{jazbec2020}; \citet{coronado2022}; \citet{coronado2023}; \citet{cruzmontana2025}; \citet{galdi2018} \\
U6b Notas pre-apertura & \citet{glasserman2025}; \citet{knuteson2020} \\
U7 y U7b Estudios de eventos & \citet{mackinlay1997}; \citet{brown1985}; \citet{kargarzadeh2026}; \citet{ardia2022}; \citet{guo2017} \\
U6 y U8 Pruebas múltiples & \citet{harvey2016}; \citet{harvey2020}; \citet{nikolopoulos2026} \\
Trabajo futuro & \citet{yang2024}; \citet{wu2026} \\
\end{longtable}
}
